\section{Prime bounds and comparison lengths}\label{adp:numerical}

For a row inducing outside $\Theta$, every character obtained by
multiplying it by a component of $\nu_i$ is nonprincipal. Otherwise
the inducing row character would belong to $\Theta$. This remains
true for each recursive child. The global hypothesis and
\cite[Lemma 4.9]{OAI26} therefore permit the prime contour to move to
$\Rea s=s_\kappa+e$, for fixed $e>0$, uniformly in the conductor
and height. The von Mangoldt formulation, removal of higher prime
powers and central normalization give
\begin{equation}\label{adp:prime}
 |Q_\psi|^2\ll Z^{\kappa z+\varepsilon_1}
       p_J(\boldsymbol W)^b(1+|\boldsymbol\omega|)^h .
\end{equation}
Here $p_J$ is a finite smooth seminorm, and $\boldsymbol\omega$
records the independent normalized twists of the profiles.
The contour displacement costs $2ez$, so it is controlled by the
bounded \emph{total} length, not by the number of factors.
The orders $J,b,h$ are fixed after the profiles and requested
derivatives. They do not add a power of $Z$.

Extra masks are removed by the exact identities
\begin{align*}
 S_{k,\mathfrak R}(X)
 &=\sum_{d\mid\mathfrak R}\mu(d)\psi_k^0(d)q_d^{-1/2}
                    T_k(X/q_d),\\
 Q_{k,i,\mathfrak R}
 &=Q_{k,i}^{\,0}-P_i^{-1/2}
       \sum_{p\mid\mathfrak R}\psi_k^0(p)\nu_i(p)W_i(q_p/P_i),
       \qquad P_i=Z^{z_i}.
\end{align*}
Here $\psi_k^0$ retains the natural zeros but omits the additional
mask, and $T_k$ is its centrally normalized plain sum.
Their absolute coefficient masses are $Z^{\varepsilon_1}$ by
\cite[Lemma 4.10 and Section 18.1]{OAI26}.
The inducing character is unchanged; each deletion decreases the
left side of \eqref{adp:admission}. If no prime factor survives,
one uses the unrestricted zero-factor assertion.

Set $A=n_1+n_2+z$. At the small-width endpoint, admission gives
\begin{equation}\label{adp:terminal}
 z\le\frac{M}{6\kappa}\le\frac{3M}{13},\qquad \kappa z\le M/6.
\end{equation}
The second inequality is the required terminal exponent payment.
The slightly larger length range in the first inequality has no
analytic consequence.

For the comparison step take $L=M/4$. If $A>5M/6$, then
\[
 (6\kappa-1)z\le M-A,\qquad 6\kappa-1\ge10/3,
 \quad\text{hence}\quad z<M/20 .
\]
If one original plain factor has length below $L$, reflect the other.
Otherwise compare the original product with a product having lengths
$L,A-z-L$, the same profiles, and the same product of scales; reflect
its longer factor when estimating the comparison.

The primitive functional equation, with the redundant deleted Euler
factors restored, gives a reflected length at most $M-n+\xi$ for
a factor of length $n$. This is exactly the reflection of
\cite[Section 18.1]{OAI26}: its conductor times its redundant radical
is bounded by $Z^M$; retained subunit annuli are rescaled only after
the support test, and both the annular tail and scale supremum are
paid by finite seminorms and polynomial heights.
Consequently the comparison has
\begin{align}
 A_{\rm comp}&\le3M/2-A+2z+\xi\le23M/30+\xi,\label{adp:comparison-a}\\
 A_{\rm comp}+(6\kappa-1)z
 &\le5M/2-2A+2z+\xi\le14M/15+\xi.\label{adp:comparison-admit}
\end{align}
Both required boundaries have a margin of $M/15$ before $\xi$.
Thus the comparison uses only the uncentered range already proved
at that width.

For completeness, suppose an aggregate clipping error $e_c$ changes
the inequality to $A\ge5M/6-e_c$. Then
$z\le M/20+3e_c/10$, and the two bounds become
\begin{equation}\label{adp:clipping}
 A_{\rm ref}\le23M/30+\xi+\frac85e_c,\qquad
 A_{\rm ref}+(6\kappa-1)z
       \le14M/15+\xi+\frac{13}{5}e_c .
\end{equation}
The coefficients $8/5$ and $13/5$ include the perturbation of $A$
itself, not only that of $z$. The zero branch of the reflected
positive part has affine bound at most
$7M/15+\xi+13e_c/10$. The long comparison factor has length at least
$8M/15-13e_c/10$. If $M\ge\rho$ and $\xi,e_c\le\rho/100$, the
margins in the uncentered boundary, admission boundary, and the
long-factor boundary $M/2$ are respectively
\[
 \frac{61\rho}{1500},\qquad
 \frac{23\rho}{750},\qquad
 \frac{61\rho}{3000}.
\]
These are positive uniformly throughout the new interval.
